\section{Future Directions}
The ARS framework establishes a hardware-conscious spatial classification architecture for parallel and streaming data processing.
While the present work focuses primarily on the Aero and Exp D implementations, several promising research directions remain open.

\subsection{Recursive Cache-Aware Refinement}
The experiments presented in this paper indicate that cache residency plays a critical role in determining refinement efficiency.
As dataset sizes continue to grow, partition sizes may eventually exceed cache capacity, exposing additional memory latency and reducing throughput.

Exp A explored recursive spatial refinement as a mechanism for preserving cache-local execution by recursively repartitioning oversized bins.
Future work may investigate adaptive cache-aware recursion strategies, dynamic refinement thresholds, and formal cache-complexity models capable of predicting optimal partition hierarchies across diverse hardware platforms.

\subsection{Hierarchical Memory Staging}
The scatter phase remains fundamentally constrained by memory-system behavior.
Exp B demonstrated that software write-combining buffers can significantly improve memory-bus utilization by transforming fragmented writes into contiguous cache-line bursts.

Future research may explore multi-level staging architectures that explicitly leverage L1, L2, and LLC hierarchies, as well as NUMA-aware staging strategies designed for high-core-count processors and multi-socket systems.

\subsection{Entropy-Adaptive Execution}
One of the most promising exploratory directions is the development of workload-adaptive execution policies.
Exp C demonstrated that runtime entropy estimation can be used to dynamically select different execution strategies depending on dataset characteristics.

Future systems may extend this concept by incorporating online learning, adaptive cost models, or lightweight predictive heuristics capable of selecting partitioning factors, refinement strategies, and mapping functions automatically at runtime.
Such mechanisms would transform ARS from a static architecture into a self-optimizing data-processing framework.

\subsection{Distributed Spatial Classification}
The current implementation targets shared-memory multi-core systems.
However, the spatial partitioning model naturally exposes coarse-grained data parallelism that may be suitable for distributed environments.

Future work may investigate cluster-scale variants of ARS in which spatial bins are distributed across nodes and refined independently.
Such architectures could potentially combine local cache-efficient classification with distributed merge-free aggregation strategies.

\subsection{Accelerator-Aware Implementations}
The ARS execution model consists primarily of histogram generation, prefix-sum computation, and spatial scattering, all of which map naturally to massively parallel hardware.

Future implementations may investigate GPU, FPGA, and heterogeneous CPU-GPU variants of the framework.
In particular, the division-free classification model employed by Aero appears well suited to SIMD and SIMT execution environments where predictable control flow is essential for maintaining high throughput.

\subsection{Persistent Streaming Architectures}
Exp D demonstrates that spatial classification can support high-throughput streaming ingestion through epoch-based micro-batching and concurrent spatial consolidation.
However, the current design remains an in-memory architecture.

Future research may investigate persistent spatial registries, fault-tolerant consolidation pipelines, and disk-backed run management systems.
Such extensions could enable ARS to serve as the foundation for large-scale streaming databases, event-processing systems, and real-time analytical platforms.

\subsection{Learned Spatial Mapping Functions}
Aero currently utilizes a sampled quantile lookup table to approximate distribution-aware partitioning.
An alternative direction involves replacing static lookup tables with lightweight learned models---drawing inspiration from learned database indexes like the Recursive Model Index (RMI) \cite{kraska2018case} and learned sorting algorithms like LearnedSort \cite{khalifa2020learnedsort}---capable of predicting optimal spatial mappings directly from observed data characteristics.

The integration of learned mapping functions may further improve occupancy balance, reduce sampling overhead, and enable adaptive behavior under rapidly changing distributions while preserving the constant-time classification properties of the framework.

Overall, the exploratory architectures and extensions presented throughout this work suggest that ARS should be viewed not as a single sorting algorithm, but as a broader family of hardware-conscious spatial processing architectures.
The Aero and Exp D implementations represent only a subset of this design space, leaving substantial opportunities for future investigation across parallel, distributed, and streaming data-processing environments.
